%%%%%%%%%%%%%%%%
%%% num 21 %%%
%%%%%%%%%%%%%%%%

\Chapter{21}

\Verse{1}
{
	\hJ{וַיִּשְׁמַע הַכְּנַעֲנִי מֶלֶךְ עֲרָד יֹשֵׁב הַנֶּגֶב כִּי בָּא יִשְׂרָאֵל דֶּרֶךְ הָאֲתָרִים וַיִּלָּחֶם בְּיִשְׂרָאֵל וַיִּשְׁבְּ מִמֶּנּוּ שֶׁבִי׃}
}
{
	\eJ{
		And the Canaanite, the king of Arad, who lived in the
		Negeb, heard that Israel was coming by the way of Atharim,
		and he fought against Israel and took some of them
		prisoners.
	}
}
{
}

\Verse{2}
{
	\hJ{וַיִּדַּר יִשְׂרָאֵל נֶדֶר לַיהֹוָה וַיֹּאמַר אִם נָתֹן תִּתֵּן אֶת הָעָם הַזֶּה בְּיָדִי וְהַחֲרַמְתִּי אֶת עָרֵיהֶם׃}
}
{
	\eJ{
		And Israel made a vow to YHWH and said, “If you will
		deliver this people into my hand, then I shall completely
		destroy their cities.”
	}
}
{
}

\Verse{3}
{
	\hJ{וַיִּשְׁמַע יְהֹוָה בְּקוֹל יִשְׂרָאֵל וַיִּתֵּן אֶת הַכְּנַעֲנִי וַיַּחֲרֵם אֶתְהֶם וְאֶת עָרֵיהֶם וַיִּקְרָא שֵׁם הַמָּקוֹם חרְמָה׃}
}
{
	\eJ{
		And YHWH listened to Israel’s voice and delivered the
		Canaanite, and they completely destroyed them and their
		cities. And the name of the place was called Hormah.
	}
}
{
}

\Verse{4}
{
	\hJ{וַיִּסְעוּ מֵהֹר הָהָר דֶּרֶךְ יַם סוּף לִסְבֹב אֶת אֶרֶץ אֱדוֹם וַתִּקְצַר נֶפֶשׁ הָעָם בַּדָּרֶךְ׃}
}
{
	\eJ{
		And they traveled from Mount Hor by way of the Red Sea
		road to go around the land of Edom. And the people’s soul
		was getting short on the way.
	}
}
{
%	\footnote{
%		the Red Sea. Hebrew yam sûp. People argue that the sea that splits in
%		Exodus 14 is not the Red Sea but rather the “Reed Sea,” an otherwise
%		unknown body of water. But the reference to the yam sûp here, when the
%		Israelites are no longer anywhere near Egypt, must refer to the eastern
%		arm of the Red Sea, which is the only body of water that extends both up
%		into Egypt and to a location far away in the Sinai. See my comment on
%		Exod 13:8. Footnote 21:4. the people’s soul was getting short. The image
%		is of getting short of breath or short of temper. The word here stands
%		in contrast to its last occurrence in the Torah: when the people cried
%		for meat, and Moses doubted whether so much meat could be provided, God
%		said, “Is YHWH’s hand getting short?!” (Num 11:23). Now they are crying
%		for food again, and it is the people’s soul or breath that is “getting
%		short.” And all of this is doubly ironic because back in Egypt the
%		people were unable to respond to Moses’ announcement of their coming
%		freedom owing to “shortage of spirit” (—Exod 6:9). Freedom has
%		not produced a speedy transformation of their nature. When faced with
%		stress they think of returning to Egypt, which now looks attractive even
%		though their condition there was described with the same term.
%	}
}

\Verse{5}
{
	\hJ{וַיְדַבֵּר הָעָם בֵּאלֹהִים וּבְמֹשֶׁה לָמָה הֶעֱלִיתֻנוּ מִמִּצְרַיִם לָמוּת בַּמִּדְבָּר כִּי אֵין לֶחֶם וְאֵין מַיִם וְנַפְשֵׁנוּ קָצָה בַּלֶּחֶם הַקְּלֹקֵל׃}
}
{
	\eJ{
		And the people spoke against {\God} and against Moses: “Why
		did you bring us up from Egypt to die in the wilderness?
		Because there’s no bread, and there’s no water, and our
		soul is disgusted with the cursed bread.”
	}
}
{
%	\footnote{
%		the people spoke against God and against Moses. The notable new element
%		in their complaint this time is that they finally seem to have
%		assimilated the fact that it is not merely Moses but God who is
%		directing them. What made them finally see this? It certainly appears to
%		be the aftermath of the events at Meribah. When Moses stopped saying
%		“It’s God, not I,” and just one time said instead “Shall we bring water
%		out,” he paid such a terrible price that the people have now seen the
%		proof that it is God, and not just Moses, who controls their fortunes.
%		Footnote 21:5. Why did you bring us up. The word “you” here is plural,
%		further reflecting that the people are, for once, not directing their
%		entire complaint against Moses. They are at least beginning to
%		acknowledge the divine power in what happened in Egypt. Even though it
%		comes in the context of a complaint, this is still a step for the
%		people, who have, all along, been unable to accept the terrifying wonder
%		of the closeness of divinity.
%	}
}

\Verse{6}
{
	\hJ{וַיְשַׁלַּח יְהֹוָה בָּעָם אֵת הַנְּחָשִׁים הַשְּׂרָפִים וַיְנַשְּׁכוּ אֶת הָעָם וַיָּמת עַם רָב מִיִּשְׂרָאֵל׃}
}
{
	\eJ{
		And YHWH let fiery snakes go among the people, and they
		bit the people, and a great many people from Israel died.
	}
}
{
}

\Verse{7}
{
	\hJ{וַיָּבֹא הָעָם אֶל מֹשֶׁה וַיֹּאמְרוּ חָטָאנוּ כִּי דִבַּרְנוּ בַיהֹוָה וָבָךְ הִתְפַּלֵּל אֶל יְהֹוָה וְיָסֵר מֵעָלֵינוּ אֶת הַנָּחָשׁ וַיִּתְפַּלֵּל מֹשֶׁה בְּעַד הָעָם׃}
}
{
	\eJ{
		And the people came to Moses, and they said, “We’ve
		sinned, because we spoke against YHWH and against you.
		Pray to YHWH, that He will turn the snake away from us.”
		And Moses prayed on behalf of the people.
	}
}
{
}

\Verse{8}
{
	\hJ{וַיֹּאמֶר יְהֹוָה אֶל מֹשֶׁה עֲשֵׂה לְךָ שָׂרָף וְשִׂים אֹתוֹ עַל נֵס וְהָיָה כּל הַנָּשׁוּךְ וְרָאָה אֹתוֹ וָחָי׃}
}
{
	\eJ{
		And YHWH said to Moses, “Make a fiery one and set it on a
		pole, and it will be that everyone who is bitten and sees
		it will live.”
	}
}
{
}

\Verse{9}
{
	\hJ{וַיַּעַשׂ מֹשֶׁה נְחַשׁ נְחֹשֶׁת וַיְשִׂמֵהוּ עַל הַנֵּס וְהָיָה אִם נָשַׁךְ הַנָּחָשׁ אֶת אִישׁ וְהִבִּיט אֶל נְחַשׁ הַנְּחֹשֶׁת וָחָי׃}
}
{
	\eJ{
		And Moses made a bronze snake and set it on a pole, and it
		was, if a snake bit a man, then he would look at the
		bronze snake and live.
	}
}
{
%	\footnote{
%		a bronze snake. When someone is bitten by a snake, the victim can look
%		at the bronze snake of Moses and be healed. This story is the only case
%		of what appears to be sympathetic magic in these books. It is also the
%		only depiction of a sacred object that is not associated with the Tent
%		of Meeting and its contents. Related to this fact, and perhaps more
%		significant, it is the only sacred object that is not connected in any
%		way with Aaron and the priests of his family. Rather, it is personally
%		made and mounted by Moses himself. This fits with the fact that it was
%		Moses’ staff that miraculously became a snake at the burning bush and in
%		Egypt (Exod 4:2–3); Aaron’s staff became a different kind of serpent
%		(Exod 7:8–10; see the comment on 7:9). The text is unclear as to whether
%		the snake is used only on the occasion of this crisis or if it continues
%		to be mounted in some visible location. On the evidence of an event
%		reported in the book of 2 Kings (18:4), it seems likely that we should
%		understand it as continuing to be in public view even after the crisis
%		is past. In that account, we are told that Moses’ snake was called
%		Nehushtan and that King Hezekiah destroyed it because the Israelites had
%		been burning incense to it. (Concerning the reason why a Judean king
%		would have dared to destroy something that was believed to have been
%		made by Moses himself, see my Who Wrote the Bible? pp. 126, 210–211.)
%		Another point: The two signs that God gives Moses at the burning bush
%		are that his staff becomes a snake and his hand becomes leprous like
%		snow (Exod 4:1–8). These now are seen to foreshadow events in Numbers:
%		the snake on a pole (Num 21:5–9) and Miriam’s becoming leprous like snow
%		(Numbers 12).
%	}
}

\Verse{10}
{
	\hE{וַיִּסְעוּ בְּנֵי יִשְׂרָאֵל וַיַּחֲנוּ בְּאֹבֹת׃}
}
{
	\eE{
		And the children of Israel traveled, and they camped in
		Oboth.
	}
}
{
}

\Verse{11}
{
	\hE{וַיִּסְעוּ מֵאֹבֹת וַיַּחֲנוּ בְּעִיֵּי הָעֲבָרִים בַּמִּדְבָּר אֲשֶׁר עַל פְּנֵי מוֹאָב מִמִּזְרַח הַשָּׁמֶשׁ׃}
}
{
	\eE{
		And they traveled from Oboth, and they camped in
		Iye-abarim, in the wilderness that is toward Moab, toward
		the sun’s rising.
	}
}
{
}

\Verse{12}
{
	\hE{מִשָּׁם נָסָעוּ וַיַּחֲנוּ בְּנַחַל זָרֶד׃}
}
{
	\eE{
		From there they traveled, and they camped in the Wadi
		Zered.
	}
}
{
}

\Verse{13}
{
	\hE{מִשָּׁם נָסָעוּ וַיַּחֲנוּ מֵעֵבֶר אַרְנוֹן אֲשֶׁר בַּמִּדְבָּר הַיֹּצֵא מִגְּבֻל הָאֱמֹרִי כִּי אַרְנוֹן גְּבוּל מוֹאָב בֵּין מוֹאָב וּבֵין הָאֱמֹרִי׃}
}
{
	\eE{
		From there they traveled, and they camped across the
		Arnon, which is in the wilderness that extends from the
		border of the Amorite, because Arnon is Moab’s border,
		between Moab and the Amorite.
	}
}
{
}

\Verse{14}
{
	\hProto{עַל כֵּן יֵאָמַר בְּסֵפֶר מִלְחֲמֹת יְהֹוָה אֶת וָהֵב בְּסוּפָה וְאֶת הַנְּחָלִים אַרְנוֹן׃}
}
{
	\eProto{
		On account of this it is said in the Scroll of the Wars of
		YHWH: Waheb in Suphah and the wadis of Arnon,
	}
}
{
%	\footnote{
%		the Scroll of the Wars of YHWH. The Torah here indicates that there are
%		known older, written sources, which it may sometimes use or cite.
%	}
}

\Verse{15}
{
	\hProto{וְאֶשֶׁד הַנְּחָלִים אֲשֶׁר נָטָה לְשֶׁבֶת עָר וְנִשְׁעַן לִגְבוּל מוֹאָב׃}
}
{
	\eProto{
		and the slope of the wadis that reached to the settlement
		of Ar, and pressed to Moab’s border.
	}
}
{
}

\Verse{16}
{
	\hE{וּמִשָּׁם בְּאֵרָה הִוא הַבְּאֵר אֲשֶׁר אָמַר יְהֹוָה לְמֹשֶׁה אֱסֹף אֶת הָעָם וְאֶתְּנָה לָהֶם מָיִם׃}
}
{
	\eE{
		And from there to Beer: that is the well of which YHWH
		said to Moses, “Gather the people, so I may give them
		water.”
	}
}
{
}

\Verse{17}
{
	\hProto{אָז יָשִׁיר יִשְׂרָאֵל אֶת הַשִּׁירָה הַזֹּאת עֲלִי בְאֵר עֱנוּ לָהּ׃}
}
{
	\eProto{
		Then Israel sang this song: Spring up, well. Sing to it.
	}
}
{
%	\footnote{
%		Then Israel sang. The poems come relatively late in the book (with the
%		possible exception of Num 12:6–8), and their effect is likely to be
%		heightened for many readers on account of this. After twenty-and-a-half
%		chapters of rebellion episodes, lists, and laws, we read that “Then
%		Israel sang” (21:17). It is the first report of song since the account
%		of the crossing of the Red Sea, which culminated with the Song of the
%		Sea (introduced with comparable words: “Then Moses and the children of
%		Israel sang,” Exod 15:1).
%	}
}

\Verse{18}
{
	\hProto{בְּאֵר חֲפָרוּהָ שָׂרִים כָּרוּהָ נְדִיבֵי הָעָם בִּמְחֹקֵק בְּמִשְׁעֲנֹתָם וּמִמִּדְבָּר מַתָּנָה׃}
}
{
	\eProto{
		The commanders hewed the well, the people’s nobles dug it,
		with a scepter, with their staffs. And from the wilderness
		to Mattanah,
	}
}
{
}

\Verse{19}
{
	\hProto{וּמִמַּתָּנָה נַחֲלִיאֵל וּמִנַּחֲלִיאֵל בָּמוֹת׃}
}
{
	\eProto{
		and from Mattanah to Nahaliel, and from Nahaliel to
		Bamoth,
	}
}
{
}

\Verse{20}
{
	\hProto{וּמִבָּמוֹת הַגַּיְא אֲשֶׁר בִּשְׂדֵה מוֹאָב רֹאשׁ הַפִּסְגָּה וְנִשְׁקָפָה עַל פְּנֵי הַיְשִׁימֹן׃}
}
{
	\eProto{
		and from Bamoth in the valley, that is in the field of
		Moab, to the top of Pisgah, which looks toward Jeshimon.
	}
}
{
}

\Verse{21}
{
	\hE{וַיִּשְׁלַח יִשְׂרָאֵל מַלְאָכִים אֶל סִיחֹן מֶלֶךְ הָאֱמֹרִי לֵאמֹר׃}
}
{
	\eE{
		And Israel sent messengers to Sihon, king of the Amorites,
		saying,
	}
}
{
}

\Verse{22}
{
	\hE{אֶעְבְּרָה בְאַרְצֶךָ לֹא נִטֶּה בְּשָׂדֶה וּבְכֶרֶם לֹא נִשְׁתֶּה מֵי בְאֵר בְּדֶרֶךְ הַמֶּלֶךְ נֵלֵךְ עַד אֲשֶׁר נַעֲבֹר גְּבֻלֶךָ׃}
}
{
	\eE{
		“Let me pass through your land. We won’t turn in at a
		field and at a vineyard and won’t drink well waters. We’ll
		go by the king’s road until we pass your border.”
	}
}
{
}

\Verse{23}
{
	\hE{וְלֹא נָתַן סִיחֹן אֶת יִשְׂרָאֵל עֲבֹר בִּגְבֻלוֹ וַיֶּאֱסֹף סִיחֹן אֶת כּל עַמּוֹ וַיֵּצֵא לִקְרַאת יִשְׂרָאֵל הַמִּדְבָּרָה וַיָּבֹא יָהְצָה וַיִּלָּחֶם בְּיִשְׂרָאֵל׃}
}
{
	\eE{
		And Sihon did not allow Israel to pass within his border.
		And Sihon gathered all his people and went out to Israel
		at the wilderness and came to Jahaz and fought against
		Israel.
	}
}
{
}

\Verse{24}
{
	\hE{וַיַּכֵּהוּ יִשְׂרָאֵל לְפִי חָרֶב וַיִּירַשׁ אֶת אַרְצוֹ מֵאַרְנֹן עַד יַבֹּק עַד בְּנֵי עַמּוֹן כִּי עַז גְּבוּל בְּנֵי עַמּוֹן׃}
}
{
	\eE{
		And Israel struck him by the sword and took possession of
		his land from Arnon to Jabbok, to the children of Ammon,
		because the border of the children of Ammon was strong.
	}
}
{
}

\Verse{25}
{
	\hE{וַיִּקַּח יִשְׂרָאֵל אֵת כּל הֶעָרִים הָאֵלֶּה וַיֵּשֶׁב יִשְׂרָאֵל בְּכל עָרֵי הָאֱמֹרִי בְּחֶשְׁבּוֹן וּבְכל בְּנֹתֶיהָ׃}
}
{
	\eE{
		And Israel took all of these cities, and Israel lived in
		all the Amorite cities in Heshbon and all of its environs,
	}
}
{
}

\Verse{26}
{
	\hE{כִּי חֶשְׁבּוֹן עִיר סִיחֹן מֶלֶךְ הָאֱמֹרִי הִוא וְהוּא נִלְחַם בְּמֶלֶךְ מוֹאָב הָרִאשׁוֹן וַיִּקַּח אֶת כּל אַרְצוֹ מִיָּדוֹ עַד אַרְנֹן׃}
}
{
	\eE{
		because Heshbon was the city of Sihon, the king of the
		Amorite, and he fought against the former king of Moab and
		took all of his land from his hand as far as Arnon.
	}
}
{
}

\Verse{27}
{
	\hProto{עַל כֵּן יֹאמְרוּ הַמֹּשְׁלִים בֹּאוּ חֶשְׁבּוֹן תִּבָּנֶה וְתִכּוֹנֵן עִיר סִיחוֹן׃}
}
{
	\eProto{
		On account of this they say proverbially: Come to Heshbon.
		Built and founded is Sihon’s city.
	}
}
{
}

\Verse{28}
{
	\hProto{כִּי אֵשׁ יָצְאָה מֵחֶשְׁבּוֹן לֶהָבָה מִקִּרְיַת סִיחֹן אָכְלָה עָר מוֹאָב בַּעֲלֵי בָּמוֹת אַרְנֹן׃}
}
{
	\eProto{
		For fire went out of Heshbon, flame from Sihon’s town; it
		consumed Ar of Moab, the lords of the high places of
		Arnon.
	}
}
{
%	\footnote{
%		high places. High places are large flat platforms where sacrificial
%		worship takes place. They are used among the Canaanites; and later the
%		Israelites use them, sometimes for worship of YHWH and sometimes for
%		worship of pagan gods. High places are forbidden to Israelites, whether
%		for YHWH or for pagan gods. Israelites are allowed to sacrifice in only
%		one place: at the Tent of Meeting, “the place where YHWH tents His name”
%		(Lev 17:5; Deut 12:5–6).
%	}
}

\Verse{29}
{
	\hProto{אוֹי לְךָ מוֹאָב אָבַדְתָּ עַם כְּמוֹשׁ נָתַן בָּנָיו פְּלֵיטִם וּבְנֹתָיו בַּשְּׁבִית לְמֶלֶךְ אֱמֹרִי סִיחוֹן׃}
}
{
	\eProto{
		Woe unto you, Moab: you’ve perished, Chemosh’s people.
		He’s made his sons refugees and his daughters captive to
		the king of the Amorite, Sihon.
	}
}
{
}

\Verse{30}
{
	\hProto{וַנִּירָם אָבַד חֶשְׁבּוֹן עַד דִּיבֹן וַנַּשִּׁים עַד נֹפַח אֲשֶׁר עַד מֵידְבָא׃}
}
{
	\eProto{
		And their fiefdom perished, Heshbon to Dibon, and we
		devastated up to Nophah which reaches to Medeba.
	}
}
{
}

\Verse{31}
{
	\hE{וַיֵּשֶׁב יִשְׂרָאֵל בְּאֶרֶץ הָאֱמֹרִי׃}
}
{
	\eE{And Israel lived in the Amorite’s land. }
}
{
}

\Verse{32}
{
	\hE{וַיִּשְׁלַח מֹשֶׁה לְרַגֵּל אֶת יַעְזֵר וַיִּלְכְּדוּ בְּנֹתֶיהָ (ויירש) [וַיּוֹרֶשׁ] אֶת הָאֱמֹרִי אֲשֶׁר שָׁם׃}
}
{
	\eE{
		And Moses sent to spy on Jazer, and they captured its
		environs and dispossessed the Amorite who was there.
	}
}
{
}

\Verse{33}
{
	\hE{וַיִּפְנוּ וַיַּעֲלוּ דֶּרֶךְ הַבָּשָׁן וַיֵּצֵא עוֹג מֶלֶךְ הַבָּשָׁן לִקְרָאתָם הוּא וְכל עַמּוֹ לַמִּלְחָמָה אֶדְרֶעִי׃}
}
{
	\eE{
		And they turned and went up the road of Bashan, and Og,
		the king of Bashan, went out to them, he and all his
		people, to war at Edrei.
	}
}
{
}

\Verse{34}
{
	\hE{וַיֹּאמֶר יְהֹוָה אֶל מֹשֶׁה אַל תִּירָא אֹתוֹ כִּי בְיָדְךָ נָתַתִּי אֹתוֹ וְאֶת כּל עַמּוֹ וְאֶת אַרְצוֹ וְעָשִׂיתָ לּוֹ כַּאֲשֶׁר עָשִׂיתָ לְסִיחֹן מֶלֶךְ הָאֱמֹרִי אֲשֶׁר יוֹשֵׁב בְּחֶשְׁבּוֹן׃}
}
{
	\eE{
		And YHWH said to Moses, “Don’t fear him, because I’ve
		delivered him and all his people and his land into your
		hand, and you’ll do to him as you did to Sihon, king of
		the Amorite, who lived in Heshbon.”
	}
}
{
}

\Verse{35}
{
	\hE{וַיַּכּוּ אֹתוֹ וְאֶת בָּנָיו וְאֶת כּל עַמּוֹ עַד בִּלְתִּי הִשְׁאִיר לוֹ שָׂרִיד וַיִּירְשׁוּ אֶת אַרְצוֹ׃}
}
{
	\eE{
		And they struck him and his sons and all his people until
		he did not have a remnant left, and they took possession
		of his land.
	}
}
{
%	\footnote{
%		they took possession of his land. The Israelites’ eastern route next
%		takes them to the territory of the Amorite kings Sihon and Og, and
%		Israel makes the same request for peaceful passage that it has made to
%		Edom. Both Sihon and Og attack Israel militarily, and both are utterly
%		defeated. Thus Israel comes to acquire territory on the eastern side of
%		the Jordan River, outside the borders of the promised land (Num
%		21:21–35).
%	}
}
